\SourceLeaf{59}$t$ congruentiae satisfaciet, id quod quaerebatur. Totum hoc artificium in eo versatur, ut numerus eruatur qui ipsius $t$, quem ignoramus, vice fungi possit. Attamen probe meminisse oportet, nos quando $\frac{p-1}{n}$ ad $n$ non est primus, supposuisse conditionem art. praec. locum habere, quae si deficit omnes conclusiones erroneae erunt; atque si regulas datas temere sequendo pro $z$ valor invenitur, cuius potestas $n^{\mathrm{ta}}$ ipsi $A$ non sit congrua. indicio hoc est, conditionem deficere adeoque methodum hanc omnino adhiberi non posse.

\Article{68}

Sed in hocce etiam casu saepe prodesse potest, hunc laborem suscepisse; operaeque pretium est, quomodo hic valor falsus ad veros sese habeat investigare. Supponamus itaque numeros $k,z$ rite esse determinatos sed $z^n$ non esse $\equiv A\,(\operatorname{mod.}p)$. Tum si modo valores expressionis ${}^{n}\!\surd{\frac{A}{z^n}}\,(\operatorname{mod.}p)$ determinari possint, hos singulos per $z$ multiplicando valores ipsius ${}^{n}\!\surd{A}$ obtinebimus. Si enim $v$ est valor aliquis ipsius ${}^{n}\!\surd{\frac{A}{z^n}}$: erit $(vz)^n\equiv A$. Sed expressio ${}^{n}\!\surd{\frac{A}{z^n}}$ eatenus hac ${}^{n}\!\surd{A}$ simplicior, quod $\frac{A}{z^n}\,(\operatorname{mod.}p)$ ad exponentem minorem plerumque pertinet quam $A$. Scilicet si numerorum $t,q$ divisor communis maximus est $d$, $\frac{A}{z^n}\,(\operatorname{mod.}p)$ ad exponentem $d$ pertinebit, id quod ita demonstratur. Substituto pro $z$ valore, fit $\frac{A}{z^n}\equiv\frac{1}{A^{kn-1}}\,(\operatorname{mod.}p)$. At $kn-1$ per $\frac{p-1}{nq}$ divisibilis (art. praec.), $\frac{p-1}{n}$ vero per $t$ (ibid.) sive $\frac{p-1}{nd}$ per $\frac{t}{d}$. Atqui $\frac{t}{d}$ ad $\frac{q}{d}$ est primus (\emph{hyp.}), quare etiam $\frac{p-1}{nd}$ per $\frac{tq}{dd}$ sive $\frac{p-1}{nq}$ per $\frac{t}{d}$, adeoque etiam $kn-1{}^{\prime}$ per $\frac{t}{d}$ et $(kn-1)d$ per $t$ erit divisibilis. Hinc $A^{(kn-1)d}\equiv1\,(\operatorname{mod.}p)$. Unde facile deducitur. $\frac{A}{z^n}$ ad potestatem $d^{\mathrm{tam}}$ evectum unitati congruum fieri. Quod vero $\frac{A}{z^n}$ ad exponentem minorem quam $d$ pertinere non possit facile quidem demonstrari potest, sed quoniam ad finem nostrum non requiritur, huic rei non immoramur. Certi igitur esse possumus, $\frac{A}{z^n}\,(\operatorname{mod.}p)$ semper ad minorem exponentem pertinere quam $A$, unico excepto casu, scilicet quando $t$ ipsum $q$ metitur, adeoque $d=t$

Sed quid iuvat, quod $\frac{A}{z^n}$ ad minorem exponentem pertinet quam $A$? Plures numeri dantur qui possunt esse $A$ quam qui possunt esse $\frac{A}{z^n}$, et quando secundum eundem modulum plures huiusmodi expressiones ${}^{n}\!\surd{A}$ evolvere occasio est, id lucramur ut plures ex eodem fonte haurire possimus. Ita \emph{ex. gr.} semper unicum saltem valorem expressionis ${}^{2}\!\surd{A}\,(\operatorname{mod.}29)$ determinare in potestate erit,
\SourceLeaf{60}si modo expressionis ${}^{2}\!\surd{-1}\,(\operatorname{mod.}29)$ valores (qui sunt $\pm12$) innotuerint. Facile enim ex art. praec. perspicitur, huiusmodi expressionum unum valorem semper directe determinari posse, quando $t$ impar, et $d$ fieri $=2$ quando $t$ par; praeter $-1$ autem nullus numerus ad exponentem $2$ pertinet.

\emph{Exempla.} Quaeritur ${}^{3}\!\surd{31}\,(\operatorname{mod.}37)$. Hic $p-1=36$, $n=3$, $\frac{p-1}{3}=12$, adeoque $q=3$: debet igitur esse $3k\equiv1\,(\operatorname{mod.}4)$ quod obtinetur ponendo $k=3$. Hinc $z\equiv31^3\,(\operatorname{mod.}37)\equiv6$, inveniturque revera $6^3\equiv31\,(\operatorname{mod.}37)$. Si valores expressionis ${}^{3}\!\surd{1}\,(\operatorname{mod.}37)$ sunt noti, etiam reliqui expr. ${}^{3}\!\surd{6}$ valores determinari possunt. Sunt vero illi $1,10,26$, per quos multiplicando ipsum $6$, prodeunt reliqui $\equiv23$ et $8$.

Si autem quaeritur valor expr. ${}^{2}\!\surd{3}\,(\operatorname{mod.}37)$, erit $n=2$, $\frac{p-1}{n}=18$; adeoque $q=2$. Hinc debet esse $2k\equiv1\,(\operatorname{mod.}9)$, unde fit $k\equiv5\,(\operatorname{mod.}9)$. Quare $z\equiv3^5\equiv21\,(\operatorname{mod.}37)$; at $21^2$ non $\equiv3$, sed $\equiv34$; est autem $\frac{3}{34}\,(\operatorname{mod.}37)\equiv-1$, atque ${}^{2}\!\surd{-1}\,(\operatorname{mod.}37)\equiv\pm6$; unde obtinentur valores veri $\pm6\,.\,21\equiv\pm15$.

Haec fere sunt, quae hic de talium expressionum evolutione tradere licuit. Palam est methodos directas satis prolixas saepe evasuras: at hoc incommodum tantum non omnibus methodis directis in numerorum theoria incumbit: neque ideo negligendum censuimus, quantum hic praestare valeant ostendere. Etiam hic observare convenit, artificia particularia quae exercitato haud raro se offerunt sigillatim explicare, non esse instituti nostri.

\begin{center}\emph{Nexus indicum in systematibus diversis.}\end{center}

\Article{69}

Revertimur nunc ad radices quas diximus primitivas. Ostendimus, radice primitiva quacunque pro basi assumta omnes numeros, quorum indices ad $p-1$ primi, etiam fore radices primitivas, nullosque praeter hos: unde simul radicum primitivarum multitudo sponte innotescit. V. art. 53. Quamnam autem radicem primitivam pro basi adoptare velimus, in genere arbitrio nostro relinquitur; unde intelligitur, etiam hic, ut in calculo logarithmico, plura quasi systemata dari posse\begingroup\renewcommand{\thefootnote}{*)}\footnote{In eo autem differunt, quod in logarithmis systematum numerus est infinitus, hic vero tantus, quantus numerus radicum primitivarum. Manifesto enim bases congruae idem systema generant.}\endgroup,
\SourceLeaf{61}quae quo vinculo connexa sint videamus. Sint $a,b$ duae radices primitivae, aliusque numerus $m$, atque, quando $a$ pro basi assumitur, index numeri $b$ $\equiv\beta$, numeri $m$ vero index $\equiv\mu\,(\operatorname{mod.}p-1)$; quando autem $b$ pro basi assumitur. index numeri $a\equiv\alpha$, numeri $m$ vero $\equiv\nu\,(\operatorname{mod.}p-1)$. Tum erit $\alpha\beta\equiv1\,(\operatorname{mod.}p-1)$; namque $a^\beta\equiv b$, quare $a^{\alpha\beta}\equiv b^\alpha\equiv a\,(\operatorname{mod.}p)$, (\emph{hyp.}), hinc $\alpha\beta\equiv1\,(\operatorname{mod.}p-1)$. Per simile ratiocinium invenitur $\nu\equiv\alpha\mu$, atque $\mu\equiv\beta\nu\,(\operatorname{mod.}p-1)$. Si igitur tabula indicum pro basi $a$ constructa habetur, facile in aliam converti potest, ubi $b$ basis. Si enim pro basi $a$ ipsius $b$ index est $\equiv\beta$, pro basi $b$ ipsius $a$ index erit $\equiv\frac{1}{\beta}\,(\operatorname{mod.}p-1)$, multiplicandoque per hunc numerum omnes tabellae indices, habebuntur omnes indices pro basi $b$.

\Article{70}

Quamvis autem plures indices numero dato contingere possint. aliis aliisque radicibus primitivis pro basi acceptis, omnes tamen in eo convenient, quod omnes eundem divisorem maximum cum $p-1$ communem habebunt. Si enim pro basi $a$, index numeri dati est $m$, pro basi $b$ vero $n$, atque divisores maximi his cum $p-1$ communes $\mu,\nu$ supponuntur esse inaequales, alter erit maior, \emph{ex. gr.} $\mu>\nu$, adeoque $\mu$ ipsum $n$ non metietur. At designato indice ipsius $a$, quando $b$ pro basi assumitur, per $\alpha$, erit (art. praec.) $n\equiv\alpha m\,(\operatorname{mod.}p-1)$ adeoque $\mu$ etiam ipsum $n$ metietur. \emph{Q. E. A.}

Hunc divisorem maximum indicibus numeri dati, ipsique $p-1$ communem, a basi non pendere, etiam inde perspicuum, quod aequalis est ipsi $\frac{p-1}{t}$, designante $t$ exponentem ad quem numerus. de cuius indicibus agitur, pertinet. Si enim index pro basi quacunque est $k$, erit $t$ minimus numerus per quem $k$ multiplicatus ipsius $p-1$ multiplum evadit (excepta cifra) vid. artt. 48, 58, sive minimus valor expressionis $\frac{0}{k}\,(\operatorname{mod.}p-1)$ praeter cifram; hunc autem aequalem esse divisori maximo communi numerorum $k$ et $p-1$ ex art. 29 nullo negotio derivatur.

\Article{71}

Porro facile demonstratur, basin ita semper accipere licere, ut numerus ad exponentem $t$ pertinens indicem quemlibet datum nansiscatur, cuius quidem maximus divisor cum $p-1$ communis $=\frac{p-1}{t}$. Designemus hunc brevitatis gratia per $d$, sitque index propositus $\equiv dm$, numerique propositi. quando quaelibet
\SourceLeaf{62}radix primitiva $a$ pro basi accipitur, index $\equiv dn$, eruntque $m,n$ ad $\frac{p-1}{d}$ sive ad $t$ primi. Tum si $\varepsilon$ est valor expressionis $\frac{dn}{dm}\,(\operatorname{mod.}p-1)$, simulque ad $p-1$ primus, erit $a^\varepsilon$ radix primitiva, qua pro basi accepta numerus propositus indicem $dm$ adipiscetur (erit enim $a^{\varepsilon dm}\equiv a^{dn}\equiv$ numero proposito), id quod desiderabatur. Sed expressionem $\frac{dn}{dm}\,(\operatorname{mod.}p-1)$ valores ad $p-1$ primos admittere, ita probatur. Aequivalet illa expressio huic: $\frac{n}{m}\,(\operatorname{mod.}\frac{p-1}{d})$ sive $\frac{n}{m}\,(\operatorname{mod.}t)$ vid. art. 31, 2, eruntque omnes eius valores ad $t$ primi; si enim aliquis valor $e$ divisorem cum $t$ communem haberet, hic divisor etiam ipsum $me$ metiri deberet, adeoque etiam ipsum $n$, cui $me$ secundum $t$ congruus, contra hypoth., ex qua $n$ ad $t$ primus. Quando igitur omnes divisores primi ipsius $p-1$ etiam ipsum $t$ metiuntur, \emph{omnes} expr. $\frac{n}{m}\,(\operatorname{mod.}t)$ valores ad $p-1$ primi erunt multitudoque eorum $=d$; quando autem $p-1$ alios adhuc divisores primos, $f,g,h$ etc. implicat, ipsum $t$ non metientes, ponatur valor quicunque expr. $\frac{n}{m}\,(\operatorname{mod.}t)\equiv e$. Tum autem quia omnes $t,f,g,h$ etc. inter se primi, inveniri potest numerus $\varepsilon$, qui secundum $t$ ipsi $e$, secundum $f,g,h$ etc. vero numeris quibuscunque ad hos respective primis fiat congruus (art. 32). Talis itaque numerus per nullum factorem primum ipsius $p-1$ divisibilis adeoque ad $p-1$ primus erit, uti desiderabatur. Tandem haud difficile ex combinationum theoria deducitur, talium valorum multitudinem fore $=\frac{p-1\,.\,f-1\,.\,g-1\,.\,h-1.\ \mathrm{etc.}}{t\,.\,f\,.\,g\,.\,h\,.\,\mathrm{etc.}}$; sed ne digressio haec in nimiam molem excrescat, demonstrationem, quum ad institutum nostrum non sit adeo necessaria, omittimus.

\begin{center}\emph{Bases usibus peculiaribus accommodatae.}\end{center}

\Article{72}

Quamvis in genere prorsus arbitrarium sit, quaenam radix primitiva pro basi adoptetur, interdum tamen bases aliae prae aliis commoda quaedam peculiaria praebere possunt. In tabula I semper numerum $10$ pro basi assumsimus, quando fuit radix primitiva; alioquin basin ita semper determinavimus ut numeri $10$ index evaserit quam minimus, \emph{i. e.} $=\frac{p-1}{t}$ denotante $t$ exponentem ad quem $10$ pertinuit. Quid vero hinc lucremur, in Sect. VI ostendemus ubi eadem tabula ad alios adhuc usus adhibebitur. Sed quoniam etiam hic aliquid arbitrarii remanere potest, ut ex art. praec. apparet: ut aliquid certi statueremus, ex omnibus radicibus primitivis quaesitum praestantibus \emph{minimam} semper pro basi elegimus. Ita pro $p=73$, ubi $t=8$ atque $d=9$, $a^\varepsilon$ habet $\frac{72\,.\,2}{8\,.\,3}$ \emph{i. e.} $6$ valores, qui sunt $5,14,20,28,39,40$. Assumsimus itaque minimum $5$ pro basi.
